%% Journal of Open Research Software Latex template

\documentclass{jors}

% Added for citations
\usepackage{natbib}
\usepackage{graphicx}

%% Set the header information
\pagestyle{fancy}
\definecolor{mygray}{gray}{0.6}
\renewcommand\headrule{}
\rhead{\footnotesize 3}
\rhead{\textcolor{gray}{JORS software Latex paper template version 0.2}}

\begin{document}

\rule{\textwidth}{1pt}

\section*{(1) Overview}

\vspace{0.5cm}

\section*{Title}

The title of the software paper should focus on the software, e.g. “Text mining software that accomplishes X”. The title should be factual, relating to the functionality of the software and the area it relates to rather than making claims about the software, e.g. “easy-to-use”. Other than perhaps the software name, the title should not contain any acronyms.

\section*{Paper Authors}

1. Last name, first name; \\
2. Last name, first name

\section*{Paper Author Roles and Affiliations}

1. First author role and affiliation \\
2. Second author role and affiliation

\section*{Abstract}

A short (approximately 100 word) summary of the software.

\section*{Keywords}

keyword 1; keyword 2; etc. Keywords should make it easy to identify who and what the software will be useful for.

\section*{Introduction}

Provide a brief introduction to the software and its purpose. Include relevant background information, and any relevant citations. A short comparison with software that implements similar functionality should be included in this section. 

This software builds upon prior work in the field. For example,  \citet{smith2020example} describe an early implementation of similar  functionality. Foundational background concepts are discussed in \citep{johnson2018book}. A related conference implementation was  introduced by \citep{williams2019conf}. As shown in Figure~\ref{fig:architecture}, the system consists of three modules.

% Although often not necessary, a diagram can be used to describe the software's architecture or other characteristics.
%Image credit: CS Odessa, CC BY 4.0 <https://creativecommons.org/licenses/by/4.0>, via Wikimedia Commons
\begin{figure}[htbp]
    \centering
    \includegraphics[width=0.8\textwidth]{diagram.png}
    \caption{System architecture diagram.}
    \label{fig:architecture}
\end{figure}

\section*{Implementation and architecture}

Describe how the software is implemented, with details of its architecture so that those unfamiliar with the software can understand how it works. Use of relevant diagrams is appropriate. Please also describe any variants and associated implementation differences. As shown in Table~\ref{tab:software_comparison}, DataStreamR provides broader feature coverage than competing tools.

% Although often not necessary, a table can be used to describe the software's features.
\begin{table}[htbp]
	\centering
	\caption{Comparison of DataStreamR with Alternative Software Tools}
	\label{tab:software_comparison}
	\begin{tabular}{lcccc}
		\hline
		\textbf{Feature} & \textbf{DataStreamR} & \textbf{Tool A} & \textbf{Tool B} & \textbf{Tool C} \\
		\hline
		Open Source & Yes & Yes & No & Yes \\
		Real-Time Processing & Yes & Limited & Yes & No \\
		Built-in Visualization & Yes & No & Yes & Limited \\
		Modular Architecture & Yes & No & Limited & Yes \\
		API Integration & Yes & Yes & Limited & No \\
		Cross-Platform Support & Yes & Yes & Yes & Limited \\
		\hline
	\end{tabular}
\end{table}

\section*{Quality control}

Detail the level of testing that has been carried out on the code (e.g. unit, functional, load etc.), and in which environments. If not already included in the software documentation, provide details of how a user could quickly understand whether the software is working (e.g. providing examples of running the software with sample input and output data). 

\section*{(2) Availability}
\vspace{0.5cm}

\section*{Operating system}

Please include minimum version compatibility.

\section*{Programming language}

Please include minimum version compatibility.

\section*{Additional system requirements}

E.g. memory, disk space, processor, input devices, output devices.

\section*{Dependencies}

E.g. libraries, frameworks.

\section*{List of contributors}

Please list anyone who helped to create the software (who may also not be an author of this paper), including their roles and affiliations.

\section*{Software location:}

{\bf Archive}

\begin{description}[noitemsep,topsep=0pt]
	\item[Name:] The name of the archive.
	\item[Persistent identifier:] DOI or similar.
	\item[Licence:] Open license.
	\item[Publisher:] Archive name.
	\item[Version published:] Version number.
	\item[Date published:] Publication date.
\end{description}

{\bf Code repository}

\begin{description}[noitemsep,topsep=0pt]
	\item[Name:] Repository name.
	\item[Identifier:] Repository URL.
	\item[Licence:] Open license.
	\item[Date published:] Publication date.
\end{description}

\section*{Language}

Language of repository, software and supporting files.

\section*{(3) Reuse potential}

Please describe in as much detail as possible the ways in which the software could be reused by other researchers both within and outside of your field. This should include use cases for the software and details of how the software might be modified or extended (including how contributors should contact you) if appropriate. Also you must include details of what support mechanisms are in place for this software.

\section*{Acknowledgements}

Please add any relevant acknowledgements to anyone else who supported the project in which the software was created, but did not work directly on the software itself.

\section*{Funding statement}

If the software resulted from funded research please give the funder and grant number.

\section*{Competing interests}

If any of the authors have any competing interests, these must be declared. The authors’ initials should be used to denote differing competing interests. For example: “BH has minority shares in [company name], which partly funded the research grant for this project. No other authors have competing interests. If there are no competing interests, please add the statement: “The authors declare that they have no competing interests.”

% Bibliography using external .bib file
\bibliographystyle{plainnat}
\bibliography{references}

\vspace{2cm}

\rule{\textwidth}{1pt}

{ \bf Copyright Notice} \\
Authors who publish with this journal agree to the following terms: \\

Authors retain copyright and grant the journal right of first publication with the work simultaneously licensed under a \href{http://creativecommons.org/licenses/by/3.0/}{Creative Commons Attribution License} that allows others to share the work with an acknowledgement of the work's authorship and initial publication in this journal. \\

Authors are able to enter into separate, additional contractual arrangements for the non-exclusive distribution of the journal's published version of the work (e.g., post it to an institutional repository or publish it in a book), with an acknowledgement of its initial publication in this journal. \\

By submitting this paper you agree to the terms of this Copyright Notice, which will apply to this submission if and when it is published by this journal.

\end{document}